\documentclass[10pt,a4paper]{amsart}
\usepackage[margin=19mm]{geometry}
\usepackage{amsmath,amssymb,mathtools}
\usepackage[scr=rsfs,cal=euler]{mathalfa}
\usepackage{xfrac}
\usepackage[unicode,hidelinks,bookmarks=false]{hyperref}
\newcommand{\ie}{i.e.\ }
\newcommand{\cf}{cf.\ }
\newcommand{\resp}{respectively\ }
\newcommand{\Subsection}{\subsection}
\providecommand{\nobreakdash}{\nobreak\hbox{-}}
\setlength{\emergencystretch}{2em}
\multlinegap0pt
\usepackage{bm}
\usepackage{bbm}
\usepackage{dsfont}
\usepackage{xspace}
\usepackage{cancel}
\usepackage{mathtools}
\usepackage{amsthm}
\makeatletter
\@ifundefined{theorem}{\newtheorem{theorem}{Theorem}}{}
\makeatother
\title[A complete classification of endomorphisms of Kiselman's sem]{A complete classification of endomorphisms of Kiselman's semigroup}
\author{Luka Andrensek}
\date{Source: arXiv:2506.03604v2 (CC BY 4.0)}
\begin{document}
\maketitle

\noindent\emph{Source and licence.} A complete classification of endomorphisms of Kiselman's semigroup. Version arXiv:2506.03604v2 (CC BY 4.0), distributed under the Creative Commons Attribution 4.0 International licence, as recorded in the arXiv licence metadata for this exact version and shown on the abstract page https://arxiv.org/abs/2506.03604v2. Full rights notice: licenses/arxiv-2506.03604-CC-BY-4.0.txt.\par\smallskip

\noindent\textit{Editorial scope.} Counted unit: the Theorem whose source label is \texttt{thm:Psi} in the article (source0.tex lines 398--404, source label \texttt{thm:Psi}), delivered together with the article's own complete original proof of it (source0.tex lines 406--475, one \texttt{proof} environment of 70 lines). No mathematics is written, repaired or re-derived: statement and proof are the article's own text line for line. Required context retained uncounted: none. Every definition, display and label of those lines is retained unchanged. Omitted from this package and not redistributed: 1--397, 405--405, 476--792. Nothing inside a retained excerpt is omitted or altered: every retained body is delivered exactly as the article's lines are written, so no omission receipt of any kind accompanies this package. Citations: the retained excerpts carry no citation command at all, so no bibliography entry of the article is transcribed and no reference is redistributed. Reference adaptation: none was needed and none was made; every reference inside the delivered unit resolves inside it, so no unresolved reference or citation is produced. Printed slips kept literal: no printed word, symbol, hypothesis or number of the counted statement or of its proof was altered. Layout and class: the article's own class is \texttt{sn-jnl} with packages including graphicx, fontenc, multirow, amsmath, amssymb, amsfonts, amsthm, mathrsfs, appendix, xcolor, textcomp, manyfoot, booktabs, algorithm; the wrapper is the lane's pinned \texttt{amsart} editorial wrapper at 10pt on A4, which restates the theorem-like environments the retained text needs and adds the article's own macro definitions that the retained text needs (none), with extra wrapper packages (bm, bbm, dsfont, xspace, cancel, mathtools). The delivered numbering is the unit's own. Assets: no retained span contains a figure environment, a graphics-inclusion command, a PDF-page inclusion, a table or a TikZ picture; no image file is copied into this package, the package declares no asset dependency and no image file is redistributed with it. Licence: version arXiv:2506.03604v2 of this article is distributed under the Creative Commons Attribution 4.0 International licence, as recorded in the arXiv licence metadata for this exact version and shown on the abstract page https://arxiv.org/abs/2506.03604v2. A full rights notice is delivered at licenses/arxiv-2506.03604-CC-BY-4.0.txt. Characters: every retained excerpt is pure ASCII, so the delivered engine, XeLaTeX with its default Latin Modern text font, reports no missing character.
\begin{theorem}\label{thm:Psi}
   The map $\Psi : M_n \rightarrow D_n$, defined by
   \[
      \Psi(X_1, X_2, \dots, X_n) = (\chi(X_1), \chi(X_2), \dots, \chi(X_n))
   \]
   where the right-hand side is interpreted as a matrix whose $i$-th column is $\chi(X_i)$, is an isomorphism of monoids $(M_n, \ast)$ and $(D_n, \cdot)$.
\end{theorem}

\begin{proof}
   We first show that $\Psi$ maps into $D_n$. Let $\mathbf{X} = (X_1, X_2, \dots, X_n) \in M_n$ and denote $M = \Psi(\mathbf{X})$.
   Observe that $M_{x, i} = 1$ if and only if $(\chi(X_i))_x = 1$, which is equivalent to $x \in X_i$.
   Now assume $i < j$ and assume that the following equalities hold
   \begin{align*}
      &M_{x, i} = 0, \quad M_{x, j} = 1, \\
      &M_{y, i} = 1, \quad M_{y, j} = 0.
   \end{align*}
   We aim to show that $x < y$ is impossible. 
   The upper equalities are equivalent to 
   \begin{align*}
      &x \notin X_i, \quad x \in X_j, \\
      &y \in X_i, \quad y \notin X_j, 
   \end{align*}
   which is further equivalent to 
   \[ 
      x \in X_j \setminus X_i \quad \text{and} \quad y \in X_i \setminus X_j. 
   \]
   We apply monotonicity of $(X_1, X_2, \dots, X_n)$ and get $x>y$. Hence, $\Psi(\mathbf{X}) \in D_n$.



   The map $\Psi$ is obviously injective.

   We now prove surjectivity. Let $M \in D_n$. Define $X_i = \{x \in \{1, 2, \dots, n\} \mid M_{x, i} = 1\}$. First, we check that $(X_1, X_2, \dots, X_n)$ is monotone.
   Let $i < j$ and suppose $x \in X_j \setminus X_i$ and $y \in X_i \setminus X_j$. Then we have
   \begin{align*}
      M_{x, i} = 0, \quad M_{x, j} = 1, \\
      M_{y, i} = 1, \quad M_{y, j} = 0.
   \end{align*}
   Since $M \in D_n$, it must follow that $x>y$, and hence $(X_1, X_2, \dots, X_n)$ is monotone.
   Furthermore, for $x, i \in \{1, 2, \dots, n\}$ we have that $(\Psi(X_1, X_2, \dots, X_n))_{x, i} = 1$ if and only if $x \in X_i$, which is equivalent to  $M_{x, i} = 1$ by the definition of $X_i$.
   Hence
   \[
   \Psi(X_1, X_2, \dots, X_n) = M,
   \]
   and surjectivity is proven. 


   Now we show that $\Psi$ is a homomorphism. Let $\mathbf{X}, \mathbf{Y} \in M_n$, where
   $\mathbf{X} = (X_1, X_2, \dots, X_n)$ and $\mathbf{Y} = (Y_1, Y_2, \dots, Y_n)$. We compute
   \begin{align*}
      \Psi(\mathbf{X} \ast \mathbf{Y}) = \Psi(\bigcup_{j \in Y_1} X_j, \bigcup_{j \in Y_2} X_j, \dots, \bigcup_{j \in Y_n} X_j).
   \end{align*}
   It follows that for $x, i \in \{1, 2, \dots, n\}$ we have 
   \begin{align*}
      \Psi(\mathbf{X} \ast \mathbf{Y})_{x, i} = 1 &\Leftrightarrow x \in \bigcup_{j\in Y_i}X_j \\
      &\Leftrightarrow \text{there exists } k \in Y_i \text{ such that } x \in X_k.
   \end{align*}
   On the other hand,
   \begin{align*}
      (\Psi(\mathbf{X}) \cdot \Psi(\mathbf{Y}))_{x, i} = \bigvee_{k=1}^n (\Psi(\mathbf{X})_{x, k} \land \Psi(\mathbf{Y})_{k, i}).
   \end{align*}
   Hence, we have
   \begin{align*}
      (\Psi(\mathbf{X}) \cdot \Psi(\mathbf{Y}))_{x, i} = 1 &\Leftrightarrow \text{there exists } k \in \{1, 2, \dots, n\} \text{ such that } \Psi(\mathbf{X})_{x, k} = \Psi(\mathbf{Y})_{k, i} = 1 \\
      &\Leftrightarrow \text{there exists } k \in \{1, 2, \dots, n\} \text{ such that } x \in X_k \text{ and } k \in Y_i \\
      &\Leftrightarrow \text{there exists } k \in Y_i \text{ such that } x \in X_k.
   \end{align*}
   We thus obtain
   \[
      \Psi(\mathbf{X} \ast \mathbf{Y})_{x, i} = (\Psi(\mathbf{X}) \cdot \Psi(\mathbf{Y}))_{x, i}.
   \]
   Since $x$ and $i$ were chosen arbitrarily, we conclude that
   \begin{align*}
      \Psi(\mathbf{X} \ast \mathbf{Y}) = \Psi(\mathbf{X}) \cdot \Psi(\mathbf{Y}).
   \end{align*}
   Furthermore, since $\Psi(\{1\}, \{2\}, \dots, \{n\}) = I$, it follows that $\Psi$ is an isomorphism of monoids. This concludes the proof.

\end{proof}

\end{document}
